\section{Introdução}

\subsection{Motivação e contexto}

A digitalização automática de partidas de xadrez a partir de imagens tem aplicações
práticas relevantes: transmissão de torneios (conversão de uma partida física em
notação digital, em tempo real, sem sensores no tabuleiro), alimentação de
\textit{engines} de análise a partir da posição capturada por câmera, acessibilidade
(descrição automática de uma partida para pessoas com deficiência visual), e como
estudo de caso para o ensino de Visão Computacional --- o tabuleiro combina geometria
regular e previsível com uma quantidade significativa de variabilidade visual (ângulo de
câmera, iluminação, material das peças).

Este projeto foi desenvolvido para a disciplina INE410121 / TRV410001 --- Visão
Computacional (UFSC) e tem como objetivo construir, do zero, uma pipeline completa de
leitura de tabuleiro: da fotografia bruta em perspectiva até a notação FEN da posição e
a identificação da jogada realizada entre dois estados consecutivos do tabuleiro.

\subsection{Por que uma abordagem híbrida (clássica + \textit{deep learning})}

A disciplina é centrada em técnicas clássicas de processamento de imagens --- domínio do
valor (histogramas, limiarização), domínio do espaço (operadores de borda, suavização,
morfologia), geometria projetiva (Hough, homografia) --- e a maior parte da pipeline
deste trabalho é resolvida com essas ferramentas. A exceção deliberada é a
\textbf{classificação do tipo de peça} (peão, torre, cavalo, bispo, dama, rei), etapa em
que abordagens clássicas (\textit{template matching}, momentos de Hu, HOG) mostraram-se
insuficientes: as doze classes de peça, fotografadas em ângulo e talhadas no mesmo
material do tabuleiro, produzem silhuetas ambíguas demais para descritores clássicos de
forma. Nessa única etapa, o projeto recorre a \textit{transfer learning} com uma
ResNet-34 pré-treinada, mantendo clássica toda a geometria (detecção de bordas, linhas,
homografia, segmentação) e a decisão de ocupação/cor.

Essa escolha não é apenas prática --- ela também é didaticamente interessante: o
resultado final expõe, lado a lado, onde a visão clássica é suficiente (leitura
geométrica do tabuleiro) e onde ela deixa de ser competitiva (reconhecimento de forma
fina sob baixo contraste), permitindo comparar diretamente o custo/benefício das duas
famílias de técnicas sobre o mesmo problema.

\subsection{Objetivos}

\begin{enumerate}[nosep]
  \item Detectar automaticamente o tabuleiro em uma foto em perspectiva e corrigir a
        distorção via homografia, obtendo uma vista de cima (\textit{top-down}).
  \item Segmentar o tabuleiro retificado em 64 casas e determinar, para cada uma, se
        está ocupada e, em caso positivo, a cor da peça.
  \item Classificar o tipo de cada peça ocupante via \textit{transfer learning}.
  \item Reconstruir a posição completa em notação FEN (todos os seis campos).
  \item Comparar dois estados consecutivos do tabuleiro e inferir a jogada realizada,
        incluindo casos especiais (captura, roque, \textit{en passant}, promoção).
  \item Quantificar o desempenho de cada etapa isoladamente e da pipeline de ponta a
        ponta, identificando o gargalo do sistema.
\end{enumerate}

\subsection{Organização do documento}

A Seção~\ref{sec:dataset} descreve o dataset utilizado. A
Seção~\ref{sec:metodologia} detalha a pipeline clássica (pré-processamento, bordas,
Hough, homografia, segmentação, ocupação e cor). A Seção~\ref{sec:deep-learning} cobre o
classificador de tipo de peça baseado em ResNet-34. A Seção~\ref{sec:fen-jogadas}
descreve a reconstrução de FEN e a detecção/notação de jogadas. A
Seção~\ref{sec:resultados} apresenta os resultados quantitativos de cada etapa e da
leitura de ponta a ponta. A Seção~\ref{sec:limitacoes} discute limitações conhecidas e
trabalhos futuros, seguida da conclusão e dos links de referência do projeto.
